% Example of a main part

\chapter{Overview}
\label{sec:grundl}
This chapter provides a quick overview for readers who do not want to read the
entire work in detail.

\chapter{Goals and Motivation}
\label{sec:real}
What should actually be achieved? Why? Why not leave everything as it is? Who
should benefit from the result?

\chapter{State of the Art}
\label{sec:real-unter}
A description of the current state of technology and science. It must also be
clearly defined how the problem relates to the system context and the known
solution space must be described (existing solutions).

\newpage
\section{Problem Description}
\label{sec:real-unterWeiter}

\newpage
\section{Existing Solutions}
\label{sec:real-unterWeiterNoch}


\section{Literature References}
\label{sec:real-literatur}

References in the text: \cite{doc:stz} and \cite{doc:gun}.

\chapter{Approach}
\label{sec:loesung}
Depending on the type of work, different solution approaches can be described
and discussed here.

\chapter{Work Carried Out}
\label{sec:erledigt}
It must be clear what was built on, i.e. what already existed, and what was
developed as part of this thesis.

\chapter{Results}
\label{sec:ergebnis}
This may contain, for example, a user guide or analysis results.
\enquote{News} measurement results.

\chapter{Summary}
\label{sec:zus}
A short summary and possibly a description of the points still missing for a
complete, market-ready solution.


